\documentclass[11pt]{article}
\usepackage[T1]{fontenc}
\usepackage{amsmath,amssymb,amsthm}
\usepackage[margin=1in]{geometry}

\newtheorem{theorem}{Theorem}
\newtheorem{lemma}{Lemma}
\newtheorem{proposition}{Proposition}
\newtheorem{remark}{Remark}

\title{Laboratory V77: Restricted Topology-Tree Transfer}
\author{NC0\_k-Avoid Laboratory}
\date{August 2026}

\begin{document}
\maketitle

\paragraph{Setting.}
Let $M$ be a set of $m$ gates with supports $\operatorname{supp}(i)$ and define
\[
 \lambda(S)=\left|\bigcup_{i\in S}\operatorname{supp}(i)
 \cap \bigcup_{i\notin S}\operatorname{supp}(i)\right|.
\]
A supplied branch decomposition is a subcubic tree $T$ whose leaves are the
gates and whose displayed cuts all have $\lambda$-value at most $b$.

\paragraph{Prior-art ingredient.}
Frederickson's restricted multilevel partitions give a binary topology-tree
hierarchy of logarithmic height. Every topology cluster is a connected vertex
set with at most three leaving tree edges, and a cluster with three leaving
edges consists of one vertex. This topology-tree framework and its relation
to top trees are prior art.

\begin{lemma}[Retained two-edge lemma]
Attach gate labels only to the degree-one leaves of $T$. Delete topology-tree
branches containing no gate label and suppress unary nodes. Every retained
label-bearing cluster has at most two source-tree edges leaving it.
\end{lemma}

\begin{proof}
A retained cluster with more than one vertex cannot have three leaving edges,
because the topology-tree invariant allows external degree three only for a
singleton. A singleton retained cluster contains a gate label and is therefore
a source-tree leaf, so it has degree one. Hence every retained cluster has at
most two leaving edges.
\end{proof}

\begin{lemma}[Two-edge support-boundary cover]
Let $C$ be a retained cluster and let $S(C)$ be its gate labels. The support
boundary of $S(C)$ is contained in the union of the middle sets of at most two
original branch-tree edges.
\end{lemma}

\begin{proof}
If a variable occurs in a gate of $S(C)$ and in a gate outside $S(C)$, the
unique path between the corresponding source leaves exits the connected
cluster $C$. It therefore crosses one of the at most two leaving edges. The
variable belongs to that edge's middle set.
\end{proof}

\begin{theorem}[Restricted topology-tree transfer]
From a supplied width-$b$ subcubic gate branch decomposition, one obtains a
rooted binary gate tree $T'$ satisfying
\[
 \operatorname{width}(T')\le 2b,
 \qquad \operatorname{height}(T')=O(\log m),
 \qquad \operatorname{EPL}(T')=O(m\log m).
\]
\end{theorem}

\begin{proof}
The prior-art topology hierarchy has height $O(\log |V(T)|)=O(\log m)$.
Pruning unlabeled branches and suppressing unary nodes preserves a rooted binary
gate hierarchy and does not increase height. At every retained node, the
preceding cover lemma and the supplied width bound give
$\lambda(S(C))\le 2b$. A tree with $m$ leaves and logarithmic height has
external path length $O(m\log m)$.
\end{proof}

\begin{proposition}[V75 consequence]
Rebuilding the V75 symbolic residual circuit on $T'$ gives incremental
prefix-avoidance time
\[
 O\!\left(m\log m\,A(2b)^2\operatorname{poly}(n,m)\right).
\]
\end{proposition}

\begin{proposition}[Cluster-level tightness]
There is a rank-three four-gate instance with supplied width $b=3$ and a valid
label-bearing two-boundary-edge cluster whose support boundary has size six.
Thus the two-edge inequality can attain $2b$ for one cluster.
\end{proposition}

\begin{remark}
This proposition is not a global factor-two lower bound: a different
logarithmic-height hierarchy may avoid that cluster. V77 does not prove that
factor two is optimal, does not refute width preservation, does not solve
unrestricted \(\mathrm{NC}^0_3\)-Avoid, and does not resolve P versus NP.
Novelty and peer review are not claimed.
\end{remark}

\paragraph{References.}
G.~N. Frederickson, \emph{Ambivalent Data Structures for Dynamic
2-Edge-Connectivity and $k$ Smallest Spanning Trees}, SIAM J. Comput. 26(2),
1997. S.~Alstrup, J.~Holm, K.~de Lichtenberg, and M.~Thorup,
\emph{Maintaining Information in Fully-Dynamic Trees with Top Trees},
arXiv:cs/0310065.

\end{document}
